\chapter{System Design}
\label{ch:design}

This chapter documents the structural decisions that shaped APEX. It walks the system from the outside in: first the deployment context that the codebase assumes, then the layered backend and the React single-page application, then the data model that everything stores against, then the API surface, and finally the migration pipeline that keeps the schema in sync with the code. The tone is descriptive rather than prescriptive --- every decision recorded here is reflected in the code that ships.

\section{Architectural overview}

APEX is a three-tier web application. A React single-page application runs in the user's browser. A single Go process serves a typed REST API on port~\texttt{8080} (\texttt{main.go:35--56}). PostgreSQL holds all durable state behind GORM~\cite{gorm}; Judge0~\cite{judge0,dosilovic2018judge0} executes student code in a sandbox; an SMTP relay carries password reset and exam reminder emails. Google Identity Services~\cite{gids} is an optional fourth dependency for OAuth sign-in.

\Cref{fig:F4.1} shows the deployment context. There is no microservice fleet: the Go binary is the only first-party service, and the rest are external dependencies. This matters because the project's working assumption is that a department or individual instructor self-hosts APEX, so every additional moving part is a tax on the operator.

\begin{figure}[H]
  \centering
  \includegraphics[width=0.95\linewidth]{F4.1_deployment-context.pdf}
  \caption{Deployment context: APEX is a single Go binary plus PostgreSQL, with Judge0, SMTP, and Google OAuth as external dependencies.}
  \label{fig:F4.1}
\end{figure}

The backend is split into layers --- transport, handlers, domain services, and persistence --- as drawn in \Cref{fig:F4.2}. Routes live in \texttt{*.RegisterRoutes} functions in each subpackage and are wired up by \texttt{main.go:42--54}. Handlers translate HTTP requests into domain calls and never speak SQL directly except through GORM. Domain services own non-trivial logic (auto-grading, reminder scheduling, finalisation arithmetic). The persistence layer is a thin wrapper over GORM (\texttt{internal/database}); model definitions sit under \texttt{internal/models}.

\begin{figure}[H]
  \centering
  \includegraphics[width=0.95\linewidth]{F4.2_layered-backend.pdf}
  \caption{Backend layers: transport, handlers, domain services, persistence.}
  \label{fig:F4.2}
\end{figure}

\Cref{fig:F4.3} shows the matching frontend structure: a \texttt{BrowserRouter}, an \texttt{AuthProvider}, and a single \texttt{AppRoutes} dispatcher in \texttt{frontend/src/app/routes.tsx}. The dispatcher mounts feature pages under \texttt{/dashboard/*} inside a shared \texttt{DashboardLayout}.

\begin{figure}[H]
  \centering
  \includegraphics[width=0.95\linewidth]{F4.3_frontend-arch.pdf}
  \caption{Frontend structure. The single SPA serves both teacher and student roles, gated by \texttt{TeacherRoute} and \texttt{StudentRoute}.}
  \label{fig:F4.3}
\end{figure}

\paragraph{Why a monolith.} A more elaborate split (separate auth service, separate grading worker, message bus) was considered and rejected. The grading workload is bursty around exam deadlines but the per-submission CPU cost is paid by Judge0, not by APEX itself; the Go process spends most of its time waiting on Judge0's HTTP response or on Postgres. Splitting the binary would buy distribution but at the cost of operational complexity that a self-hosting instructor cannot reasonably absorb. The grading "workers" are simply goroutines launched from \texttt{exam.SubmitAttempt} (\texttt{internal/exam/handler.go:615--622}); under load they queue inside Judge0, not inside Go.

\section{Data model}

The persistence schema centres on three aggregates: identity (users, profiles, password reset tokens), classes and enrolment (classes, members, announcements, exam-class links), and assessment (exams, problems, test cases, attempts, submissions, test results). Every model file under \texttt{internal/models} is a plain Go struct annotated for GORM. \Cref{fig:F4.4} flattens the lot into one entity-relationship diagram.

\begin{figure}[H]
  \centering
  \includegraphics[width=\linewidth]{F4.4_er-diagram}
  \caption{Entity-relationship diagram covering all GORM models in \texttt{internal/models/}.}
  \label{fig:F4.4}
\end{figure}

Two design choices in the data model are worth flagging.

\paragraph{Attempts vs submissions.} The codebase distinguishes the act of taking an exam (\texttt{ExamAttempt}, in \texttt{internal/models/exam\_attempt.go}) from the answers given to individual problems (\texttt{Submission}, in \texttt{internal/models/submission.go}). An attempt is one row per \texttt{(user\_id, exam\_id)} pair --- enforced at the database level by \texttt{uniqueIndex:idx\_user\_exam\_attempt} on both columns --- and groups the submissions for that sitting via \texttt{exam\_attempt\_id}. This shape is what makes resume-after-crash and re-aggregation arithmetic both possible and tractable. \Cref{fig:F4.7} is the class diagram for this aggregate.

\begin{figure}[H]
  \centering
  \includegraphics[width=0.95\linewidth]{F4.7_class-attempt-aggregate.pdf}
  \caption{Attempt aggregate: one \texttt{ExamAttempt} groups many \texttt{Submission} rows; each \texttt{Submission} groups its \texttt{TestResult}s.}
  \label{fig:F4.7}
\end{figure}

\paragraph{Polymorphic problems.} A \texttt{Problem} can be \texttt{coding}, \texttt{mcq}, or \texttt{written}. Rather than a separate table per type, fields specific to each variant live on the same row and are nil/empty for the other types: MCQ uses \texttt{Options} and \texttt{CorrectOptionIDs}; written uses \texttt{MaxWordCount}, \texttt{Rubric}, \texttt{RequireManualGrading}; coding uses \texttt{StarterCode} plus its child \texttt{TestCase} rows. \Cref{fig:F4.6} is the exam-side class diagram. The trade-off is a wider table in exchange for a much simpler grading dispatch path (\texttt{exam.SubmitAttempt} switches on \texttt{Problem.Type}) and a simpler editor on the frontend.

\begin{figure}[H]
  \centering
  \includegraphics[width=0.95\linewidth]{F4.6_class-exam-aggregate.pdf}
  \caption{Exam aggregate: \texttt{Exam}, \texttt{Problem}, \texttt{TestCase}, with the polymorphic problem fields shown explicitly.}
  \label{fig:F4.6}
\end{figure}

\Cref{tab:state-fields} summarises the lifecycle-bearing columns that drive state transitions later in this chapter.

\begin{table}[H]
  \centering
  \caption{Lifecycle-bearing columns. Each column controls a state transition documented later.}
  \label{tab:state-fields}
  \small
  \begin{tabularx}{\linewidth}{l l X}
    \toprule
    Table & Column & Role \\
    \midrule
    \texttt{exams}         & \texttt{is\_draft}                & blocks attempts; cannot be cleared without \texttt{start\_time} \\
    \texttt{exams}         & \texttt{start\_time, end\_time}   & exam window (\texttt{StartAttempt} enforces) \\
    \texttt{exams}         & \texttt{reset\_at}                & cache-bust marker for student sessions \\
    \texttt{exams}         & \texttt{reminder\_sent\_at}       & in-app reminder dedup \\
    \texttt{exams}         & \texttt{email\_reminder1h\_sent\_at} & 1-hour email dedup \\
    \texttt{exams}         & \texttt{email\_reminder\_start\_sent\_at} & "starting now" email dedup \\
    \texttt{exam\_attempts}& \texttt{status}                   & \texttt{in\_progress} $\to$ \texttt{submitted} \\
    \texttt{exam\_attempts}& \texttt{draft\_answers, draft\_saved\_at} & autosave payload, cleared on submit \\
    \texttt{exam\_attempts}& \texttt{graded\_notified}         & idempotent guard for "Exam Graded" notification \\
    \texttt{submissions}   & \texttt{status}                   & \texttt{pending} $\to$ \texttt{running} $\to$ accepted/wrong/CE/TLE/RE/pending\_review \\
    \texttt{password\_reset\_tokens} & \texttt{used\_at, expires\_at} & one-shot 1-hour reset tokens \\
    \bottomrule
  \end{tabularx}
\end{table}

\section{Domain state machines}

\Cref{fig:F4.14}, \Cref{fig:F4.15}, and \Cref{fig:F4.16} together describe the three machines that govern correctness in APEX: exams, submissions, and attempts.

\begin{figure}[H]
  \centering
  \includegraphics[width=0.95\linewidth]{F4.14_state-exam}
  \caption{Exam state machine. Publishing requires a non-null \texttt{start\_time} (\texttt{exam.UpdateExam} returns 400 otherwise, \texttt{internal/exam/handler.go:222--233}).}
  \label{fig:F4.14}
\end{figure}

\begin{figure}[H]
  \centering
  \includegraphics[width=0.95\linewidth]{F4.15_state-submission}
  \caption{Submission state machine. \texttt{pending\_review} is the special state for written answers and for MCQs that have no configured correct answers.}
  \label{fig:F4.15}
\end{figure}

\begin{figure}[H]
  \centering
  \includegraphics[width=0.95\linewidth]{F4.16_state-attempt}
  \caption{Attempt state machine. The unique \texttt{(user\_id, exam\_id)} index makes every state transition idempotent on retry.}
  \label{fig:F4.16}
\end{figure}

The publishing rule deserves a sentence of its own. Earlier revisions of the exam handler silently set \texttt{start\_time = now()} when a teacher tried to publish without a schedule. Two field-test users complained that exams opened the moment they clicked \emph{Publish}, ahead of the lesson. The handler now rejects publish-without-schedule with HTTP 400 (\texttt{internal/exam/handler.go:222--233}); the user-facing fix is one line of error-display logic in \texttt{ExamBuilder.tsx}.

\section{Sequence diagrams}

The five sequences in this section cover the operations that touch more than one subsystem: login, Google OAuth, autosave-and-submit, auto-grading of coding answers, and reminder-tick. Manual grading of written answers has a sixth sequence in \Cref{fig:F4.12}.

\begin{figure}[H]
  \centering
  \includegraphics[width=0.95\linewidth]{F4.8_seq-login}
  \caption{Login. Password verification uses bcrypt~\cite{provos1999bcrypt}; the JWT~\cite{rfc7519} is HS256 with a 24-hour expiry.}
  \label{fig:F4.8}
\end{figure}

\begin{figure}[H]
  \centering
  \includegraphics[width=0.95\linewidth]{F4.9_seq-google-oauth}
  \caption{Google OAuth. The frontend gets an ID token from Google Identity Services~\cite{gids}, the backend validates it via the official Go SDK (\texttt{idtoken.Validate}), and either matches by \texttt{google\_id}, links to an existing email, or creates a new account once a role has been chosen.}
  \label{fig:F4.9}
\end{figure}

\begin{figure}[H]
  \centering
  \includegraphics[width=\linewidth]{F4.10_seq-attempt-autosave-submit}
  \caption{Autosave + submit. The autosave endpoint (\texttt{exam.AutosaveAttempt}, \texttt{internal/exam/handler.go:650--685}) stores the raw JSON body from the client, so the SPA can serialise whatever shape its UI needs.}
  \label{fig:F4.10}
\end{figure}

\begin{figure}[H]
  \centering
  \includegraphics[width=\linewidth]{F4.11_seq-autograde-coding}
  \caption{Auto-grading. \texttt{judge0.Grade} (\texttt{internal/judge0/grader.go:13--109}) compares Judge0's normalised stdout against the test case's normalised expected output and aggregates a percentage.}
  \label{fig:F4.11}
\end{figure}

\begin{figure}[H]
  \centering
  \includegraphics[width=\linewidth]{F4.12_seq-grade-written}
  \caption{Manual grading of a written answer. \texttt{submission.GradeSubmission} (\texttt{internal/submission/handler.go:76--128}) re-checks ownership before applying the score and triggers the same finalisation pipeline as auto-grading.}
  \label{fig:F4.12}
\end{figure}

\begin{figure}[H]
  \centering
  \includegraphics[width=\linewidth]{F4.13_seq-reminder-tick}
  \caption{Reminder scheduler. \texttt{reminder.runOnce} (\texttt{internal/reminder/scheduler.go:29--53}) runs once a minute, in three sweeps: in-app reminder, "starts in 1 hour" email, and "starting now" email. Each sweep dedupes via a per-exam timestamp.}
  \label{fig:F4.13}
\end{figure}

\section{API surface}

\Cref{fig:F4.17} draws the route tree. Every route lives under \texttt{/api}; the \texttt{public} group is open, and the \texttt{protected} group is gated by \texttt{middleware.Auth}. Inside the protected group, role-restricted handlers gate themselves with \texttt{middleware.RequireRole("teacher")} or \texttt{("student")}, and resource ownership is re-checked inside the handler (e.g. exam handlers use \texttt{WHERE id = ? AND teacher\_id = ?}).

\begin{figure}[H]
  \centering
  \includegraphics[width=\linewidth]{F4.17_api-route-tree.pdf}
  \caption{API route tree. Public routes are limited to authentication and password reset; everything else requires a JWT.}
  \label{fig:F4.17}
\end{figure}

A few API decisions are worth noting up front.

\begin{itemize}
  \item Per-question submissions are not exposed. The only way for a student to create submission rows is \texttt{POST /student/exams/:id/submit}, which is server-driven (\texttt{exam.SubmitAttempt}). \texttt{POST /submissions/run} only runs sample tests; it never persists. This shape eliminates a class of "submit one answer at a time, race to a stale attempt" bugs.
  \item Autosave (\texttt{PUT /student/exams/:id/autosave}) is intentionally schemaless: the handler reads the raw body and only validates that it parses as JSON. The SPA can evolve its draft shape without a coordinated server change. The trade-off is that the server cannot grade or repair drafts; that is fine, because drafts are cleared on \texttt{Submit}.
  \item Manual grading (\texttt{PUT /submissions/:id/grade}) is a single endpoint that handles both initial scoring and override. The handler re-verifies ownership and then calls \texttt{judge0.FinalizeAttempt} so the student's aggregate score updates in lock-step.
\end{itemize}

\section{Authentication, authorisation, and roles}

Authentication is a stateless JWT~\cite{rfc7519}. \texttt{auth.generateToken} (\texttt{internal/auth/handler.go:475--485}) signs an HS256 token with a 24-hour expiry and the claims \texttt{user\_id}, \texttt{email}, \texttt{role}, \texttt{name}, \texttt{exp}. \texttt{middleware.Auth} (\texttt{internal/middleware/auth.go:11--58}) parses and validates the token, then drops the three identity claims into the Gin context with comma-ok type assertions; a token whose \texttt{user\_id} claim is anything other than a JSON number is rejected with HTTP 401. \texttt{middleware.RequireRole} (\texttt{internal/middleware/auth.go:61--81}) is a small helper that gates a route group by role string.

The \texttt{config.Load} step at startup refuses to boot if \texttt{JWT\_SECRET} is unset --- so production deployments cannot accidentally fall back to a default secret. The original code path silently used a development default if the variable was empty; that was patched in commit \texttt{2eaac88} (see \Cref{ch:security}).

OAuth follows OpenID Connect~\cite{oidc-core,rfc6749}. The SPA obtains an ID token from Google Identity Services and posts it to \texttt{POST /auth/google}. The server validates the token against \texttt{GOOGLE\_CLIENT\_ID} via the official Go \texttt{idtoken.Validate}, then either matches by \texttt{google\_id}, links to an existing email account, or creates a new user once the SPA has prompted for a role. There is no password fallback for Google-created accounts; their \texttt{password\_hash} stays empty and \texttt{auth.Login} explicitly rejects empty hashes (\texttt{internal/auth/handler.go:112--115}).

\section{Migration pipeline}

GORM's \texttt{AutoMigrate} can add columns and create tables, but it cannot run data backfills, drop tables, or change a column's nullability without losing data. APEX therefore runs a hand-written, idempotent migration step \emph{before} \texttt{AutoMigrate}. \Cref{fig:F4.18} draws the order; \texttt{internal/database/migrations.go} is the source of truth.

\begin{figure}[H]
  \centering
  \includegraphics[width=0.95\linewidth]{F4.18_migration-pipeline}
  \caption{Schema bring-up order. Idempotent SQL runs first, then GORM's \texttt{AutoMigrate} creates new tables and indices, then a small post-step adjusts foreign-key constraints that depend on tables \texttt{AutoMigrate} just created.}
  \label{fig:F4.18}
\end{figure}

Each migration block is gated on the underlying table existing (\texttt{tableExists}, \texttt{internal/database/migrations.go:8--12}). On a first-boot deployment those gates skip the blocks; \texttt{AutoMigrate} then creates the tables for real, and re-running the binary on the second boot replays the gates without effect. The most consequential block in the file is the \texttt{exam\_attempts.score} backfill at lines 176--203, which re-aggregates scores over the full \texttt{problem} list of each exam: an earlier code path silently dropped skipped questions out of the denominator, so a student who answered only the easy questions could appear at 100\,\%. The backfill is gated on all three of \texttt{exam\_attempts}, \texttt{problems}, and \texttt{submissions} existing.

The conventions documented in \texttt{ORM\_RULES.md} (idempotent SQL, gated on table existence, run before \texttt{AutoMigrate}, no destructive statements without an \texttt{IF EXISTS} clause) are the contract between code reviewers and migrations: a new column is fine in a model file, but a backfill on existing rows must be added to \texttt{RunMigrations}.

\section{Cross-cutting concerns}

\paragraph{CORS.} \texttt{middleware.CORS} (\texttt{internal/middleware/cors.go:10--18}) allows only \texttt{http://localhost:5173}, the Vite dev server. Production deployments terminate TLS and serve the SPA from the same origin as the API, so CORS is not relied on at the trust boundary. This is intentional: a permissive CORS policy is one of the standard misconfigurations behind cross-site request forgery~\cite{cwe-352}.

\paragraph{Notifications.} \texttt{notification.Create} (\texttt{internal/notification/helper.go:11--45}) inserts a row in \texttt{notifications} and, asynchronously, sends an email if the user opted in. The async email is HTML-escaped via \texttt{html.EscapeString} on both \texttt{title} and \texttt{description}; this was the fix landed for the email HTML-injection issue described in \Cref{ch:security}.

\paragraph{Reminders.} The reminder scheduler (\texttt{internal/reminder/scheduler.go}) is started from \texttt{main} as a goroutine. It sleeps for 15 seconds at boot to let the database settle and then fires every minute. Each sweep dedupes via a per-exam timestamp column (\texttt{reminder\_sent\_at}, \texttt{email\_reminder1h\_sent\_at}, \texttt{email\_reminder\_start\_sent\_at}), so a process restart inside the same minute does not cause duplicate notifications.

\paragraph{Finalisation.} The arithmetic that turns a set of submissions into an attempt score lives in one place, \texttt{judge0.maybeFinalizeAttempt} (\texttt{internal/judge0/grader.go:245--328}). Both the auto-grader and the manual-grading endpoint call it. The function iterates the exam's full problem list (so skipped questions stay in the denominator at score~0), excludes \texttt{pending\_review} submissions from both numerator and denominator until a teacher grades them, then fires a single "Exam Graded" notification once nothing is left in \texttt{pending\_review}. The \texttt{exam\_attempts.graded\_notified} flag makes that notification idempotent.

\section{Summary}

The shape of APEX is intentionally small. One Go binary, one Postgres database, one optional sandbox, one optional SMTP relay; one polymorphic problem table, one attempt aggregate that owns its submissions, one finalisation function that everyone calls. The three state machines in \Cref{fig:F4.14}, \Cref{fig:F4.15}, and \Cref{fig:F4.16} fit on three pages because the underlying objects do not have unprincipled side states. The next chapter walks the implementation that fills in this shape, and \Cref{ch:security} examines the threats that the shape is designed to resist.
